\section{Conclusion}\label{sec:conclusion}
Using part-of-speech (POS) and lemma information in subword tokenizers in five languages, we assess them under one protocol: gold morpheme boundaries, held-out text, repeated training samples and a language model probe.
Weighting subword merges by POS transitions does not change token boundaries or output, because the transition probability only describes the order of words, not internal word structure.
Weighting subword merges within content words changes the vocabulary as desired, but does not always give a better accuracy than without.

Using lemmas to deduce word boundaries instead increases boundary recall by 23 to 77 points, at the cost of 4--8\% more tokens.
How much of the gain carries over to unseen words depends on the distinctiveness of a language's affixes.
If a tokenizer is to include morphological boundaries, it should use a lemmatizer or segmentation lexicon, where a tagger does not provide any value.
However, the quality of these boundaries for language models is a question in itself.
A language model probe ranks these tokenizers differently than the boundary scores, with the unigram model first and Lemma2Word last in each language, and will need larger scale than our 7M-parameter language model to resolve.
Morphological alignment cannot yet be recommended as a goal.
